\section{Pendahuluan}

\subsection{Latar Belakang}

Layanan digital kini menjadi bagian dari cara orang merencanakan perjalanan. Pada 2020, Asia Tenggara memiliki sekitar 400 juta pengguna internet aktif dan 94\% di antaranya berniat tetap memakai layanan digital setelah pandemi \cite{econ2020}. Nilai transaksi perjalanan daring di Indonesia mencapai US\$9 miliar pada 2024, tumbuh 24\% per tahun \cite{econ2024}, sementara perjalanan wisatawan nusantara Januari--Agustus 2024 mencapai 674,6 juta, tumbuh 19,2\% \cite{bps2024}. Rekomendasi yang tepat menghemat waktu wisatawan dalam mencari dan memutuskan.

Namun, rekomendasi wisata tidak cukup hanya menebak destinasi yang disukai. Sistem harus menemukan \emph{kombinasi} destinasi beserta \emph{urutan} kunjungannya yang pas dengan anggaran, waktu, dan jumlah orang. Dari 201 destinasi terdapat $2^{201} \approx 3{,}2 \times 10^{60}$ kemungkinan himpunan destinasi, belum termasuk urutannya. Pendekatan \textit{greedy} (pilih yang termurah atau terdekat dulu) mudah terjebak pada pilihan yang tampak baik secara lokal. Algoritma Genetika (GA) \cite{holland1975, goldberg1989} cocok untuk masalah seperti ini karena menjelajahi banyak solusi sekaligus lewat populasi, diarahkan oleh fungsi fitness, dan konvergen ke solusi yang sangat baik dalam waktu singkat.

\subsection{Tujuan}

\begin{enumerate}
  \item Memodelkan rekomendasi rencana perjalanan sebagai gabungan masalah pemilihan (mirip Knapsack) dan pengurutan rute (mirip TSP), dengan batas anggaran dan waktu.
  \item Membangun GA lengkap tanpa pustaka GA siap pakai: representasi kromosom, fungsi fitness, seleksi, crossover, mutasi, elitisme, dan kriteria berhenti.
  \item Membuktikan efektivitas GA terhadap metode pembanding (dua varian greedy dan pencarian acak) serta menganalisis sensitivitas parameternya.
  \item Menyajikan hasilnya dalam aplikasi web interaktif yang dapat memakai dataset pengguna sendiri.
\end{enumerate}

\subsection{Sumber Kode}

Seluruh kode ditulis dalam TypeScript dan tersedia di \url{https://github.com/kukuhagushermawan/GA-rekomendasi}. Mesin GA berada di folder \texttt{src/ga/} dan dipisah per tahap agar mudah ditelusuri (\Tabelref{tab:kode}).

\begin{table}[H]
\centering
\caption{Pemetaan berkas kode ke isi laporan.}
\label{tab:kode}
\small
\begin{tabular}{@{}l >{\raggedright\arraybackslash}p{8.4cm} l@{}}
\toprule
\textbf{Berkas} & \textbf{Isi} & \textbf{Subbab} \\
\midrule
\texttt{data.ts}       & membaca dan memvalidasi dataset, menyusun matriks per nama & \ref{sec:data} \\
\texttt{kromosom.ts}   & representasi kromosom dua bagian dan dekodenya & \ref{sec:kromosom} \\
\texttt{fitness.ts}    & biaya, durasi, validitas, fitness, jadwal harian & \ref{sec:model}, \ref{sec:fitness} \\
\texttt{operator.ts}   & seleksi, crossover, dan mutasi & \ref{sec:operator} \\
\texttt{ga.ts}         & populasi awal, loop GA, elitisme, kriteria berhenti, top-5 & \ref{sec:operator} \\
\texttt{pembanding.ts} & greedy, pencarian acak, eksperimen & \ref{sec:eksperimen} \\
\texttt{worker.ts}     & menjalankan GA di Web Worker untuk aplikasi web & \ref{sec:implementasi} \\
\texttt{tests/ga.test.ts} & 37 uji unit & \ref{sec:verifikasi} \\
\bottomrule
\end{tabular}
\end{table}
